\documentclass[name=רשימות פונקציות ממשיות]{../../../tex/classes/noChapterMegasumPreamble}
\usetikzlibrary{calc}
\usepackage{../../../tex/packages/mathShortcuts}
\usepackage{../../../tex/packages/hebrewSupport}
\usepackage{../../../tex/packages/megasumextendedtheoremsupport}

\newcommand \lectext[1] {{\color{blue}\textbf{#1}}}
\newcommand \firstlec   {\lectext{תחילת הרצאה 1}}
\newcommand \nthlec[2]  {\lectext{סוף הרצאה #1, תחילת הרצאה #2}}

\author{שחר פרץ}
\title{\name}
\date{2026 סמסטר ב'}
\begin{document}
	\renewcommand{\footrule}{\rule{\linewidth-19pt}{0.25pt}\vspace{-5pt}}
	\firstPage{2027A $\sim$ מבוסס על הקלטות 2016B}
	\newpage

	\hfil \textit{תאריך קימפול: \today}

	\setcounter{section}{-1}
\section[מבוא למבוא]{\en{Introduction to the Introduction}}
\subsection[רישיון]{License $\sim$ רישיון}
	\en{The following summary is provided under the GNU General Public License version 3 (GPLv3). It can be distributed and/or modified under the terms of the license, or any later version of it. Additional information can be found \href{https://www.gnu.org/licenses}{here}. }

	הסיכום להלן מסופק תחת רשיון התוכנה החופשית של גנו גרסה 3 (GNU General Public License (GPL) version 3). ניתן להעתיקו ו/או להפיצו תחת GPLv3 או גרסה מאוחרת יותר. מידע נוסף אפשר למצוא \href{https://www.gnu.org/licenses}{כאן}.

\subsection{הרבה מילים שאפשר לדלג עליהן}
	להלן מובא סיכום מוקלד של ההרצאות \href{https://youtu.be/CrTJoGaCI8U?si=Ccje3Wyot-WmqLMW}{ביוטיוב} של פרופסור בועז קלרטג ב־2016. הסיכום הוקלד בסמסטר הראשון של שנת 2027. {\color{red} הסיכום עדיין בשלבי כתיבה}.


\subsection{סימונים}
	בסיכום הבא נניח את הסימונים הבאים:
	\begin{itemize}
		\item \hfil $[n] \cod \N \cap [0, n]$
		\item \hfil $\subset \; \cod \; \subseteq$
		\item \hfil $\bigcup A \cod \bigcup_{X \in A}X$
		\item \hfil $I \cod [0, 1]$
	\end{itemize}



	\newpage
	\setcounter{tocdepth}{3} % show subsubsections
	\tableofcontents

	%	\listoftheorems

	\middleText{\textit{בתיאבון}}

	\section[מבואות למיניהם]{\en{All Kinds of Introductions}}
	\firstlec
	\subsection{חומרי קריאה}
	\begin{itemize}
		\item Stein \& Sakarachi: Measure Theory
		\item לינדשטראוס, וייס: מבוא לאנליזה מודרנית 
	\end{itemize}
	
	\subsection{השלמות מטופולוגיה}
	הרציונלים בני מנייה, כלומר נוכל למספר $\{q_n\}_{n = 1}^{\infty} = \Q$. עתה נתבונן בקבוצה הבאה: 
	\[ A = \bigcup_{n = 1}^{\infty}\cl{q_n - \frac{1}{2^{n}}, q_n + \frac{1}{2^{n}}} \]
	נבחין שהיא איחוד של קבוצות ולכן פתוחה. משום ש־$\Cl \Q = \R$, מתבקש לשאול האם $A = \R$. נוכל לענות על שאלה זו באמצעות מושג ה\textit{מידה} שנגדיר בהמשך, שיגדיר את גודל הקבוצה להיות סופי, ובפרט: 
	\[ m(A) \le \sum_{n = 1}^{\infty} \frac{1}{2^{n}} = 2 \]
	
	\textit{הערה: }משפט דיריכלה על קירובים דיופאנטים: 
	\[ \forall p, q \in \N\co \sof{\sqrt 2 - \frac{p}{q}} > \frac{1}{q^{2}} \]
	
	לכן, לדוגמה, ''אי אפשר לקרב את $\frac{1}{\sqrt 2}$ בצורה יחסית טובה ב־$A$``. 
	
	שימו לב שהקטעים בהגדרה של $A$ לא זרים, בהכרח, מצפיפות. עם זאת, ניתן להציג את $A$ כאיחוד קטעים זרים ופתוחים (היא לא קשירה). 
	\theo{\label{theo:1}תהא $A \subset \R$ קבוצה פתוחה. אז אפשר להציג את $A$ כאיחוד סופי או בן־מנייה של איחוד קטעים פתוחים זרים. }
	\defi{שקטע $I \subset A$ קטע פתוח מקסימלי ביחס ל־$A$ אם לא קיים $J \subset \R$ קטע פתוח כך ש־$I \subsetneq J \subset A$}
	נוכל להגדיר את המושג באופן דומה בעבור קבוצה פתוחה. 
	\rmark{לצורך המשפט, ''קטע פתוח`` במובן הרחב, כלומר בפרט קרנות $(\infty, a), (a, \infty)$ ו־$\R$. }
	\lem{תהא $A$ פתוחה. אז לכל $x_0 \in A$ קיים קטע מקסימלי שמכיל את $x_0$. }
	\begin{proof}
		נסמן: 
		\begin{gather*}
			\bg_* = \sup\ccb{\bg > x_0 \co [x_0, \bg) \subset A} \\
			\ag_* = \inf\ccb{\ag < x_0 \co (\ag, x_0] \subset A}
		\end{gather*}
		נבחין במספר הבחנות: 
		\begin{itemize}
			\item $\bm{\ag_* < x_0 < \bg_*}$\textbf{: }כי $x_0$ נקודת פנים ($x_0 \in \Int A$ כי $A$ פתוחה) לכן $x_0 \in (\ag, \bg) \subset A$ אז יש איזשהם $\ag, \bg$ לא טרוויאלים
			\item \textbf{הקטע }$\bm{(\ag_*, \bg_*) \subset A}$\textbf{: }ניקח $x \in (\ag_*, \bg_*)$. 
			\begin{itemize}
				\item אם $x > x_0$ אז קיים $\bg \ge x_0$ כך ש־$x \in [x_0, \bg) \subset A$
				\item אם $x \le x_0$ אז $\ag_* < x$ ולכן קיימת $\ag < x$ כך ש־$x \in [\ag, x_0) \subset A$
			\end{itemize}
			\item \textbf{הקטע $I \cod \bm{(\ag_*, \bg_*)}$ מקסימלי} (ביחס ל־$A$)\textbf{: }נניח בשלילה שלא היה מקסימלי, ואז היה קטע $(\ag, \bg)$ כך ש־$(\ag_*, \bg_*) \subsetneq (\ag, \bg) \subset A$, אך מהגדרת הסופרמום $\bg \le \bg_*$ ומהאינפימום $\ag \ge \ag_*$ כלומר $(\ag, \bg) =(\ag_*, \bg_*)$ בסתירה. 
		\end{itemize}
	\end{proof}
	נפנה להוכחת המשפט. 
	\begin{proof}[הוכחת משפט \ref{theo:1}]
		 נתונה קבוצה $A$. נביט באוסף הקטעים המקסימליים ביחס ל־$A$. 
		 \begin{itemize}
		 	\item \textbf{הקטעים הנ''ל פתוחים: }טרוויאלי, קטע פתוח הוא קבוצה פתוחה. 
		 	\item \textbf{איחודם מכסה את $\bm{A}$: }ישירות מטענת העזר, סביב כל נקודה יש קטע פתוח מקסימלי כזה. 
		 	\item \textbf{הקטעים הנ''ל זרים: }נניח ש־$I, J \subset A$ שני קטעים מקסימליים שנחתכים, ובשלילה $I \cap J \neq \vn$. אזי $I \cup J \subset A$ פתוח, ומכיל את הקטע המקסימלי $I$ כלומר $I = I \cup J$. באותו האופן $J = I \cup J = I$ ולכן הם שווים. 
		 	\item \textbf{האיחוד לכל היותר בן בנייה: }כי $\R$ הוא CII – בכל קטע ישנו רציונלי, הרציונלים שווים זה מזה (הם זרים) ומאקסיומת הבחירה הבת מנייה הגדרנו העתקה חח''ע לקבוצה בת מנייה. 
		 \end{itemize}\envendproof
	\end{proof}

	\subsection{מידת לבג ב־$\R^2$}
	
	כיצד נגדיר את מושג המידה? הרעיון הוא להתאים לקבוצה $A \subset \R^{d}$ מספר, שאמור לתאר את \textit{המידה} שלה. 
	\begin{itemize}
		\item כש־$d = 1$, \textit{מידה} משמעה אורך
		\item כש־$d = 2$, \textit{מידה} משמעה אורך
		\item כש־$d = 3$, \textit{מידה} משמעה נפח
	\end{itemize}
	
	נגדיר את מידת לבג במקרה הפרטי בו $d = 2$. הבנייה עובדת באותה המידה בכל הממדים. לאחר מכן נגדיר את המידה באופן כללי בכל מרחב מידה (כלומר קבוצה עם מבנה נוסף שמתאר אילו קבוצות מדידות ואילו לא). 
	
	תהא $A \subset \R^{2}$, ונרצה באיזשהו אופן להגדיר את $m(A)$ שמייצג את השטח שלה. אנחנו כבר יודעים לחשב שטח של מלבן $a \times b$ ששטחו $a \cdot b$. זו הנחת הייסוד שנעבוד איתה (ביחד עם אדטיביות של קבוצות זרות). 
	
	\defi{\textit{מלבן אלמנטרי} ב־$\R^{2}$ הוא קבוצה מהצורה $[a, b) \times [c, d)$ כאשר $a \le b, c \le d \in \R$. }
	
	נבחין שהגדרה זו בפרט כוללת את הקבוצה הריקה. נגדיר את השטח/מידה של מלבן אלמנטרי כזה להיות: 
	\[ m(A) \cod m_2(A) \cod (b - a) \cdot (d - c) \]
	(מסמנים לעיתים ב־$m_d$ את המידה ל־$\R^{d}$). 
	 
	\begin{Lemma}[אדטיביות]
		נניח ש־$R, R_1 \dots R_N \subset \R^{2}$ מלבנים אלמנטריים, כך ש־$R = \biguplus_{k = 1}^{N}R_k$ (איחוד זר). אזי: 
		\[ m(R) = \sum_{k = 1}^{N}m_2(R_k) \]
	\end{Lemma}
	\prooftrivial
	(הטריק שעושים כאן, הוא לקחת את כל המלבנים, ''להמשיך`` את הצלעות שלהם לישרים אינסופיים, להוכיח למקרה הזה ואז לעשות רדוקציה למקרה המקורי). 
	
	באותו האופן, נוכל להוכיח את הלמה הבאה: 
	\begin{Lemma}[תת־אדטיביות]
		יהיו $P, P_1 \dots P_N \subset \R^{2}$ מלבנים אלמנטריים, כך ש־$P \subset \bigcup_{i = 1}^{N}P_i$, אזי $m(P) \le \sum_{i = 1}^{N}m(P_i)$. 
	\end{Lemma}
	
	לעיתים גם \textit{תת־אדטיביות} מכונה \textit{סאב־אדטיביות}. 
	
	נקבע מלבן אלמנטרי $\Og \subset \R^{2}$, ונעבוד בה''כ בתוכו לצורך הנוחות. 
	\defi{קבוצה $A \subset \Og$ נקראת \textit{קבוצה אלמנטרית} אם היא איחוד סופי של מלבנים אלמנטריים. }
	
	ההצגה לא יחידה (ניתן לייצגה כאיחוד סופי של מלבנים אלמנטריים שונים), אך מידת הקבוצה לא תלויה בהצגה. 
	
	\clm{ב־$\Og$, איחוד, חיתוך ומשלים של מספר סופי של קבוצות אלמנטריות, הוא קבוצה אלמנטרית. }
	\begin{proof}יהיו $P_1 \dots P_N, R_1 \dots R_T$ מלבנים אלמנטריים, כך ש־$B = \bigcup_{i = 1}^{T}R_i, \ A = \bigcup_{i = 1}^{N}P_i$ אלמנטריים. 
		\begin{itemize}
			\item איחוד: ברור ש־$A \cup B = \bigcup_{i = 1}^{T} R_i \cup \bigcup_{j = 1}^{N}P_i$ כנדרש. 
			\item חיתוך: בדומה לאיחוד, בשילוב עם העובדה שחיתוך מלבנים אלמנטריים נותנת מלבן אלמנטרי. 
			\item משלים: אם $\cg = [a, b) \times [c, d)$ אז משלים של מלבן מדה־מורגן $R = [x, y) \times [z, w) \subset \Og$ ואז נקבל ש־: 
			\[ \ol R = [a, b) \times [c, z) \cup [a, b) \times [w, c) \cup [a, x) \times [z, w) \cup [y, b) \times [z, w) \]
			קבוצה אלמנטרית. ואז ל־$A$ כללית נקבל ש־$\Og \setminus A = (\Og \setminus R_1) \cap \cdots \cap (\Og \setminus R_N)$
		\end{itemize}
	\end{proof}

	\nthlec{1}{2}
	
	\rmark{תוצאה של המשפט לעיל היא המקרה הפרטי לפיו $A \cup B, A \cap B, A \trio B, A^{c} = \Og \setminus A, A \setminus B$ כולן קבוצות אלמנטריות. מכאן ש־: 
	\begin{itemize}
		\item קבוצת הקבוצת האלמנטריות המוכלות ב־$\Og$ היא אלגברה בוליאנית. 
		\item כל $A \subset \Og$ אלמנטרית אפשר להציג כאיחוד \textit{זר} של מספר סופי של מלבנים אלמנטריים. 
	\end{itemize}}
	ההוכחה של המסקנה השנייה היא הוכחה דומה למה שכבר ראינו של להמשיך את המלבנים. 
	
	\textbf{תזכורת: }$A \trio B \cod (A \setminus B) \cup (B \setminus A)$
	
	\defi{אוסף תת־קבוצות של $\Og$ הסגור למשלים ולאיחוד סופיים, וכן מכיל את הקבוצה הריקה, נקרא \textit{אלגברה של קבוצות}. }
	\clm{\label{theo:boolring}בהינתן אלגברה של קבוצות מעל $\Og$, תחת ההגדרה: 
	\begin{align*}
		A + B \cod A \trio B && A \cdot B = A \cap B
	\end{align*}
	מקבלים חוג עם יחידה (שהוא $\vn$), שבו $\forall x \co x^{2} = x$. קבוצה שמקיימת תנאים אלו נקראת \textit{אלגברה בוליאנית}. }
	\defi{חוג שבו $\forall x \co x^{2} = x$ נקרא \textit{חוג בוליאני}. }
	\cola{חוג בוליאני היא חוג קומטטיבי ממאפיין $2$. }
	\rmark{יש איזומורפיזם בין קטגוריית החוגים הבוליאנים לבין קטגוריית האלגבראות הבוליאניות, שמנוסח במפורש במשפט \ref{theo:boolring}. לרוב משתמשים במושג אלגברה בוליאנית בהקשרים של תקב''צ, ובחוג בוליאני בהקשרים אלגבריים. }
	
	\defi{תהי $A \subset \R^{2}$ קבוצה אלמנטרית $A = \bigcup_{i = 1}^{N}P_i$ כאשר $P_i \subset \R^{2}$ מלבנים אלמנטריים. נגדיר את מידת $A$ להיות: 
	\[ m(A) \cod \sum_{i = 1}^{N}m(P_i) \]}
	\clm{\ 
	\begin{enumerate}
		\item $m(A)$ לא תלויה בהצגה של $A$ כאיחוד של מלבנים אלמנטרים זרים, 
		\item ההגדרה לעיל מזדהה עם ההגדרה הקודמת במקרה ש־$A$ מלבן אלמנטרי. 
	\end{enumerate}}
	ההוכחה של $1$ פשוטה מאדטיביות של מידת מלבנים, ו־$1$ גורר ישירות את $2$. 
	
	\clm{תכונות בסיסיות של מידת קבוצות אלמנטריות: 
	\begin{itemize}
		\item בהינתן $A, B \subset \Og$ אלמנטריות אז $A \cap B = \vn \implies m(A \cup B) = m(A) + m(B)$
		\item אם $A_1 \dots A_N$ אלמנטריות זרות בזוגות, אז $m\cl{\bigcup A_i} = \sum m(A_i)$ (באינדוקציה על $1$)
		\item אם $A_1 \dots A_N$ אלמנטריות, אז $m\cl{\bigcup A_i} \le \sum m(A_i)$ 
	\end{itemize}}
	
	ההוכחה של כולן מסתמכת על ''הרחבה`` של המלבנים. 
	
	\begin{Theorem}[תת־סיגמה־אדטיביות]\label{theo:r2subsigmaadd}
		נניח ש־$A, A_1 \dots \subset \R^{2}$ סדרת קבוצות אלמנטריות כך ש־$A \subset \bigcup_{n = 1}^{\infty} A_n$, אז: 
		\[ m(A) \le \sum_{n = 1}^{\infty}m(A_n) \]
	\end{Theorem}
	\rmark{בהקשרים הללו, האות סיגמה לרוב מסמנת איחוד בן־מנייה. }
	התכונה לעיל קשורה לשלמות של $\R^{2}$ כמרחב מטרי. אם לדוגמה היינו עוברים ברציונליים, אז לדוגמה המלבן $I = [0, 1)^{2} \cap \Q^{2}$ היה איחוד של ''המון ישרים דקים`` וזה לא היה עובד. 
	
	הוכחת המשפט מסתמכת על רדוקציה מאיחוד בן־מנייה לאיחוד סופי (שם יש לנו בבירור סאב סיגמה אדטיביות מתכונות האלגברה), והדרך המתבקשת לעשות את זה היא באמצעות היינה־בורל. 
	\lem{תהא $A$ אלמנטרית ו־$\eg > 0$. אז קיימות אלמנטריות $A', \tl A\subset \R^{2}$ כך ש־:
	\begin{enumerate}[(A)]
		\item קירוב מבפנים ע''י קבוצה סגורה: $\Cl A' \subset A \ \land \ m(A \setminus A') < \eg$. 
		\item קירוב מבחוץ ע''י קבוצה פתוחה: $A \subset \Int(\tl A) \ \land \ m(\tl A \setminus A) < \eg$. 
	\end{enumerate}}\begin{proof}
		במלבנים, ההוכחה פשוטה (בונים מלבן קטן). במקרה של איחוד סופי של $n$ מלבנים, נקרב באופן זהה כל אחד מהמלבנים בעבור $ \frac{\eg}{n}$, ואז נאחד. פנים וסגור אינווריאנטים לאיחוד ולכן סיימנו. 
	\end{proof}
	''אנחנו לא פיזיקאים`` $\sim$ המרצה אחרי שהוא כתב את ההוכחה הזו. 
	
	\begin{proof}[הוכחת \ref{theo:r2subsigmaadd}]
		לאיחוד סופי – טרוויאלי. \textbf{הרעיון: }במקום להוכיח ש־$x \le y$, נוכיח ש־$\forall \eg \co x \le y + \eg$. נקבע $\eg > 0$ כלשהו. מטרתנו להוכיח ש־$m(A) \le \sum_{n = 1}^{\infty} m(A_n)$. נחליף את $A$ באלמנטרית $\Cl A' \subset A$ כך ש־$m(A \setminus A') \le \frac{\eg}{2}$, ובאופן דומה נחליף את $A_n$ באלמנטרית $\tl A_n$ כך ש־$A \subset \Int(\tl A_n)$ ו־$m(\tl A_n \setminus A_n) < \frac{\eg}{2^{n + 1}}$. נשים לב ש־$A'$ פרה־קומפקטית (היא חסומה + היינה־בורל) ולכן: 
		\[ \Cl A' \subset A \subset \bigcup_{n = 1}^{\infty}A_n \subset \bigcup_{n = 1}^{\infty}\Int(\tl A_n) \]
		הקבוצות הפתוחות $\Int(\tl A_1) \dots \Int(\tl A_n)$ כיסוי פתוח של הקומפט $\Cl A'$, ולכן ויש תת־כיסוי סופי $\ol{A'} \subset \bigcup_{r = 1}^{\infty}$. עתה, נשתמש בטענה לאיחוד סופי של אלמנטריות: 
		\[ \begin{WithArrows}
			m(A) &= m(A') + m(A \setminus A') \\
			&\le \frac{\eg}{2} + m(A') \Arrow[ll]{$A' \subset \bigcup_{r = 1}^{N}\tl A_n$}\\
			&\le \frac{\eg}{2} + \sum_{r = 1}^{N}m(\tl A_n) \\
			&\le \frac{\eg}{2} + \sum_{n = 1}^{\infty}\csb{m(A_n) + \frac{\eg}{2^{n + 1}}} \Arrow[ll]{$\sum_{n = 1}^{\infty}\frac{1}{2^{n + 1}} = \frac{1}{2}$}\\
			&= \eg + \sum_{n = 1}^{\infty}\cl{m(A_n)}
		\end{WithArrows} \]
	\end{proof}
	
	\cola{\label{col:elemntaryprop}אז מה אנחנו יודעים על קבוצות אלמנטריות ב־$\Og$? 
		\begin{enumerate}
			\item הן מהוות אלגברה (סגורות לאיחוד סופי ומשלים)
			\item המידה $m$ אדטיבית־סופית עבורן (כלומר $A \cap B = \vn \implies m(A \cup B) = m(A) + m(B)$)
			\item $m$ היא תת־$\sg$־אדטיבית (כלומר $A \subset \bigcup_{n = 1}^{\infty}A_n \implies m(A) \le \sum_{n = 1}^{\infty}m(A_n)$)
	\end{enumerate}}

	\subsubsection{מידה חיצונית}
	
	אלו שלושת התכונות \textit{היחידות} שאנחנו עומדים להשתמש בהן בהגדרת המידה לקבוצות שאינן אלמנטריות. הניסיון הראשון שלנו להגדיר מידה כזו תהיה מידה חיצונית. 
	
	\defi{\label{defi:outermeasure}תהא $A \subset \R^{2}$ קבוצה כלשהי. אז \textit{המידה החיצונית של $A$} היא: 
	\[ m^{*}(A) = \inf \ccb{\sum_{n = 1}^{\infty}m(P_n)\ \middle\vert \ A \subset \bigcup_{n = 1}^{\infty}P_n \ \mathrm{s.t.} \ \text{$A_n$ אלמנטריים}} \]}
	הבחנה חשובה היא ש־$\Img m^{*} = [0, +\infty]$. יש קבוצות שהשטח עבורן ''אינסופי``. 
	
	\begin{itemize}
		\item \textbf{מה טוב? }אכן קבוצות שהיינו מצפים לשטח מסוים יצאו סבירות, לדוגמה השטח של מעגל היחידה הוא $\pi$. 
		\item \textbf{מה רע? }אם מאמינים באקסיומת הבחירה, ניתן לבנות $A, B \subset I^{2}$ \textit{זרות} כך ש־$m^{*}(A \cup B) < m^*(A) + m^*(B)$
	\end{itemize}
	
	הפתרון של הדבר הרע יהיה לצמצם את $m^*$ לתת קבוצה נחמדה של קבוצות ''מדידות``. 
	
	\rmark{את הרחבת מושג השטח, אפשר לבסס רק על שלוש התכונות במסקנה \ref{col:elemntaryprop}. }
	
	
	\begin{Collary}[תכונות של $m^*$]
		\begin{enumerate}
			\item אם $A \subset \R^{2}$ ברת מנייה, אז $m^*(A) = 0$. למה? כי $A = \{x_i\}_{i = 1}^{\infty}$, ואז נוכל לכסות את $x_i$ על מלבן אלמנטרי ששטחו $\frac{\eg}{2^{i}}$, ולקבל $\forall \eg > 0 \co m^*(A) \le \eg$. 
			\item בהגדרה של $m^*(A)$, אפשר להחליף ''מלבנים אלמנטריים`` ב־''קבוצות אלמנטריות`` (כי איחוד בן נייה של קבוצות סופיות הוא בן מנייה). 
			\item בהגדרה של $m^*(A)$, אפשר להחליף את $A \subset \bigcup_{n = 1}^{\infty}A_n$ בתנאי $A \subset \bigcup_{n = 1}^{\infty}\ol{A_n}$ או בתנאי $A \subset \bigcup_{n= 1}^{\infty}\Int(A_n)$
			\item עבור קבוצה $A \subset \R^{2}$ אלמנטרית, מתקיים $m(A) = m^*(A)$. 
		\end{enumerate}
	\end{Collary}
	ההבחנה השנייה, תאפשר לנו בהמשך להגדיר מידה חיצונית נטו בעזרת מעין מידה אלמנטרית (מקיימת את שלושת התכונות) מעל אלגברה נחמדה מספיק. 
	
	\begin{proof}[הוכחה של מסקנה 4]
		ברור ש־$m^*(A) \le m(A)$ בעבור כיסוי עם איבר 1 – $A$ עצמה. יותר קשה להראות ש־$m^*(A) \ge m(A)$, כלומר צ.ל. בהינתן $A \subset \bigcup_{n = 1}^{\infty}A_n$ כאשר $A_n$ אלמנטריות, מתקיים $m(A) \le \sum_{n = 1}^{\infty} m(A_n)$. גם כאן התשובה פשוטה, מתת־$\sg$־אדטיביות. 
	\end{proof}
	
	''ביטול. אין הפסקה``. 
	
	\textbf{דוגמאות של מידה חיצונית. }
	\begin{enumerate}
		\item $m^*([0, 1))^{2} = 1$ (כי זה אלמנטרי)
		\item $m^*(I^{2}) = 1$. חסם תחתון ברור מהסעיף הקודם. החסם העליון כי עבור $\eg > 0$ קל להראות ש־$m^*(A) \le 1 + \eg$ ע''י הכיסוי $[0, 1) \cup P_1 \cup P_2 \cup P_3 \cup P_4$ כאשר $P_1 \dots P_4$ מלבנים אלמנטריים סביב צלעות המלבן ממידה $\frac{\eg}{4}$ (העובי שלהם קטן כרצוננו). 
		\item באופן דומה להוכחה של $P_1 \dots P_4$, ברור ש־$m^*(I) = 0$ בתוך $\R^{2}$. 
	\end{enumerate}
	
	\theo{מידה חיצונית היא תת־$\sg$־אדטיבית, כלומר: 
	\[ m^*\!\Bigg(\smash{\overbrace{\bigcup_{n =1}^{\infty}A_n}^{A}}\Bigg) \!\le \sum_{n = 1}^{\infty}m^*(A_n) \]}\begin{proof}
		ניקח $\eg > 0$ ונוכיח את הא''ש עד לכדי אפסילון. לכל $n$, קיימות אלמנטריות $P_{n, 1} \dots P_{n, 3} \dots \subset \R^{2}$ כך ש־$A_n \subset \bigcup_{k =1}^{\infty}P_{n, k}$,		וכמעט מתקבל ה־$\inf$, כלומר: 
		\[ \begin{WithArrows}
			\sum_{k = 1}^{\infty}m^*(P_{n, k}) &< m^*(A_n) + \frac{\eg}{2^{n}} \Arrow[up]{$A \subset \bigcup_{n, k}P_{n, k}$}\\
			m^*(A) &\le \sum_{n, k = 1}^{\infty}m^*(P_{n, k}) \le \sum_{n = 1}^{\infty}m^*(A_n) + \frac{\eg}{2^{n}} = \eg + \sum_{n = 1}^{\infty}m^*(A_n)
		\end{WithArrows} \]
		כנדרש. 
	\end{proof}
	
	\subsubsection{קבוצות מדידות}
	
	בעוד כמה שבועות, נראה $A, B$ (''כלל וכלל לא אלמנטריות``) זרות כך ש־$m^*(A \cup B) < m^*(A) + m^*(B)$. זו הבעיה העיקרית במידה חיצונית – זה אמור להיות אדטיבי ביחד לשתי קבוצות זרות. 
	
	הפתרון יהיה להצטמצם לקבוצה בשם קבוצת הקבוצות המדידות, שעליהן $m^*$ אדטיבית (מידת איחוד זר היא סכום המידות) ''רוב הקבוצות הן מדידות``, בהיבט הזה שחייבים להשתמש באקסיומת הבחירה כדי לבנות קבוצות לא מדידות. 
	
	לעת עתה, אנחנו עדיין עובדים בתחום תחום קבוע $\Og \subset \R^{2}$ כלשהו. 
	
	\defi{קבוצה $A \subset \Og$ נקראת \textit{מדידה} אם ניתן לקרבה ע''י אלמנטריות, כלומר לכל $\eg > 0$ קיימת אלמנטרית $P_\eg \subset \Og$ כך ש־$m^*(A \trio P_\eg) < \eg$. }
	\rmark{בלעז, measurable. }
	\rmark{יש גישה אלטרנטיבית, שקולה, בה מגדירים את $A$ להיות מדידה אמ''מ לכל $B$ מתקיים ש־$m^*(A) = m^*(A \cup B) + m^*(A \setminus B)$. אנחנו נעבוד עם הגישה המוגדרת לעיל, שהיא הגישה של לבג. }
	''נוכיח שזה שקול?`` ''יום אחד``
	
	\textbf{דוגמאות: }\begin{enumerate}
		\item כל אלמנטרית היא מדידה, ניתנת לקירוב ע''י עצמה. 
		\item עיגולים הם גם סבבה: עם בנייה מגעילה מספיק אפשר לעשות את זה. ניתן להוכיח בצורה יותר אלגנטית באמצעות קריטריון דרבו לאינטגרביליות רימן. נראה בהמשך שכל פתוחה/סגורה היא מדידה. 
		\item אם $A \subset \Og$ קב' בעלת מידה חיצונית $0$ היא בהכרח מדידה (ניתן לקרבה עם הקבוצה הריקה). 
	\end{enumerate}
	
	בתרגיל, נוכיח ש־$d(A, B) = m^*(A \trio B)$ מקיים את א''ש המשולש וזו פסאודומטריקה. יתרה מכך, הפונקציה $A \to m^*(A)$ רציפה (מליפשיציות). 
	
	\clm{אוסף הקבוצות המדידות ב־$\Og$ סגור למשלים ולאיחוד (מהווה אלגברה). }\begin{proof}
		יהיו $A, B \subset \Og$ מדידות, ונוכיח ש־$A \cup B$ מדיד. יהי $\eg >0$. רוצים למצוא $\eg$־קירוב אלמנטרי ל־$A \cup B$. ישנן אלמנטריות $P_\eg, Q_\eg$ שמקרבות את $A, B$ בהתאמה ברזולוציה של $\frac{\eg}{2}$ (כלומר $m^*(A \trio P_\eg) < \frac{\eg}{2}$ ו־$m^*(B \trio P_\eg) < \frac{\eg}{2}$). 
		
		נרצה לקרב את $A \cup B$ באמצעות $P_\eg \cup Q_\eg$. ואכן: 
		\[ (A \cup B) \trio (P_\eg \cup Q_\eg) \subset (A \trio P_\eg) \cup (B \trio Q_\eg) \]
		ומתת־אדטיביות של $m^*$ קיבלנו ש־$m^*((A \cup B) \trio (P_\eg \cup Q_\eg)) < \frac{\eg}{2} + \frac{\eg}{2} = \eg$ וסיימנו. 
		
		מה לגבי $\Og \setminus A$? יש אפסילון־קירוב ל־$A$ ע''י $P_\eg$ כך ש־$m^*(A \trio P_\eg) < \eg$. אך $A^c \trio P_\eg^c = A \trio P_\eg$ כלומר $m^*(A^c \trio P_\eg^c) = \eg$ כנדרש. 
	\end{proof}
	
	\clm{נניח שיש $A, B \subset \Og$ \textit{מדידות} וזרות. אז $m^*(A \cup B) = m^*(A) + m^*(B)$. }\begin{proof}
		א''ש אחד ידוע מתת־אדטיביות של $m^*$. הכיוון השני (צ.ל. $m^*(A \cup B) \le m^*(A) + m^*(B)$) לאו דווקא ברור. 
		
		כדי להוכיח את הא''ש השני, ניקח $\eg > 0$ ונסתפק בלהוכיח ש־$m^*(A \cup B) \ge m^*(A) + m^*(B) - 10\eg$. ממדידות $A, B$ קיימות אלמנטריות $P_\eg, Q_\eg \subset \Og$ כך ש־$m^*(A \trio P_\eg), m^*(B \trio Q_\eg) < \eg$. 
		
		[שימו לב שלמרות ש־$A \cap B = \vn$, ייתכן ש $P_\eg \cap Q_\eg \neq \vn$, אבל אפשר להראות שהחיתוך קטן]
		
		נבחין ש־:
		\[ P_\eg \cap Q_\eg \subset (P_\eg \setminus A) \cup (Q \eg \setminus B) \subset (P_\eg \trio A) \cup (Q_\eg \trio B) \]
		לכן $m^*(P_\eg \cap Q_\eg) \le m^*(P_\eg \trio A) + m^*(Q_\eg \trio B) < 2\eg$. מכיוון ש־$P_\eg, Q_\eg \subset \Og$ אלמנטריות, אז $m^*(P_\eg) = m(P_\eg)$ וכנ''ל לגבי $Q_\eg$ (כאשר $m$ המידה שהגדרנו רק לאלמנטריות, והוכחנו בה אדטיביות סופית) ולכן: 
		\[ m^*(P_\eg \cup Q_\eg) = m^*(P_\eg) + m^*(Q_\eg) - m^*(Q_\eg \cap P_\eg) \]
		עבור $x, y \in \R$ נסמן $x \approx y$ אם $\sof{x - y} < \eg$. לעיל הסקנו ש־$\smash{m^*(P_\eg \cup Q_\eg) \overset{2\eg}{\approx} m^*(P_\eg) + m^*(Q_\eg)}$. 
		
		בנוסף, ברור ש־$m^*(A) \approx m^*(P_\eg)$ ו־$m^*(B) \approx m^*(Q_\eg)$. 
		
		לגבי $A \cup B$, ראינו בהוכחה הקודמת ש־$m^*\csb{(A \cup B) \trio (P_\eg \cup Q_\eg)} \le m^*(A \trio P_\eg) + m^*(B \trio P_\eg) \le 2\eg$. לכן $m^*(A \cup B) \overset{2\eg}{\approx} m^*(P_\eg \cup Q_\eg)$. נציב למעלה ונקבל: 
		\[ m^*(A \cup B) \overset{2\eg + 2\eg + \eg + \eg}{\approx} m^*(A) + m^*(B) \]
		ובפרט עד לכדי $10\eg$ כנדרש. 
	\end{proof}
	
	\cola{\begin{enumerate}
			\item אפשר להמשיך באינדוקציה על $A_1 \dots A_n$ מדידות וזרות בזוגות ולקבל $m^*(\bigcup A_i) = \sum m^*(A_i)$. 
			\item עבור $A, B \subset \Og$ מדידות, נוכל לקבל הכלה והפרדה: 
			\[ m^*(A \cup B) = m^*(A) + m^*(B) - m^*(A \cap B) \]
			(כי $m^*(B) = m^*(A \cap B) + m^*(B \setminus A)$ ובאופן דומה על $A$)
	\end{enumerate}}
	
	''המשפט הבא הוא פלא``. ''במידה מסוימת, הקורס הזה הוא הערות למשפט הזה``
	\begin{Theorem}[סיגמה־אדטיביות של המידה]
		תהיינה $A_1, A_2 \dots \subset \Og$ סדרה בת־מנייה של קבוצות מדידות. אז: 
		\begin{enumerate}
			\item איחוד הבן מנייה $\bigcup_{n = 1}^{\infty}A_n$ של קבוצות מדידות הוא מדיד
			\item אם הקבוצות זרות בזוגות, אז $m^*\cl{\bigcup_{n = 1}^{\infty}A_n} = \sum_{n = 1}^{\infty}m^*(A_n)$
		\end{enumerate}
	\end{Theorem}
	\begin{proof}
		\begin{itemize}
			\item להוכחת החלק הראשון של המשפט, אפשר להניח בה''כ שהן זרות בזוגות, שכן אפשר תמיד להחליפן ב־: 
			\[ \begin{cases}
				A_1' = A_1 \\
				A_n' = A_n \setminus \cl{\bigcup_{i = 1}^{n - 1} A_i}
			\end{cases} \]
			קיבלנו קבוצות $A_1' \dots A_n'$ זרות, מדידות (קבוצת המדידות היא אלגברה, ו־$A_n'$ מתקבלת ע''י מספר סופי של חיתוכים ואיחודים), ואיחודן אותו הדבר ($\bigcup A_n = \bigcup A_n'$). 
			\item לכל $N$ סופי, ידוע ש־: 
			\[ \sum_{n= 1}^{N}m^*(A_n) = m^*\cl{\bigcup^{N}_{n = 1} A_n} \le m^*(A) < \infty \]
			מתת־$\sg$־אדטיביות. לכן הטור $\sum^{\infty}_{n  =1}m^*(A_n)$ מתכנס, כי סדרת הסכומים החלקיים חסומה (הרגע הוכחנו) ומונוטונית עולה (המידה חיובית). לכן הזנב קטן כרצוננו, כלומר לכל $\eg > 0$ ישנו $N_\eg$ כך ש־$\sum_{n = N_\eg}^{\infty}m^*(A_n) < \eg$ (זנב כל טור מתכנס). 
			
			\item למה $A = \bigcup_{n = 1}^{\infty}$ מדידה? יהי $\eg > 0$, נמצא קירוב אלמנטרי מתאים. הקבוצה $\bigcup_{n= 1}^{N_\eg}A_n$ מדידה, ולכן יש אלמנטרית $P_\eg$ שמקרבת אותה $m^*(P_\eg \trio \bigcup^{N_\eg}_{n = 1}A_n) < \eg$. נראה ש־$m^*(P_\eg \trio A) < 2\eg$. אכן: 
			\[ \begin{WithArrows}
				m^*(P_\eg \trio A) &= m^*\!\cl{P_\eg \trio \csb{\bigcup_{n = 1}^{N_\eg} \cup\  \bigcup_{\mathclap{n = N_\eg + 1}}^{\infty} A_n}} \Arrow[up]{סאב־אדטיביות} \\
				&= m^*\!\cl{P_\eg \trio \bigcup_{n = 1}^{n_\eg}} + m^*\!\cl{\bigcup_{n = N_\eg + 1}^{\inft}A_n} \\
				&= \eg + \eg = 2 \eg
			\end{WithArrows} \]
			בכך הוכחנו את חלק א' של המשפט. 
			\item עתה, נותר להראות ש־$m^*\cl{\bigcup_{n = 1}^{\infty}} = \sum_{n = 1}^{\infty}m^*(A_n)$ כאשר $A_n$ זרות. כיוון אחד יש לנו מתת־$\sg$־אדטיביות. 
			
			לגבי הכיוון השני, ממונוטוניות. לכל $N$ מתקיים: 
			\[ m^*\!\cl{\bigcup_{n = 1}^{\infty}A_n} \ge m^*\!\cl{\bigcup_{n = 1}^{N}A_n} = \sum_{n = 1}^{N}m^*(A_n) \]
			ולכן מסנדוויץ' סיימנו. 
		\end{itemize}\envendproof
	\end{proof}
	
	\renewcommand{\footrule}{\rule{\linewidth-26pt}{0.25pt}\vspace{-5pt}}
	
	\defi{אוסף תת־קבוצות של $\Og$ הסגור למשלים ולאיחודים בני מנייה, וכן מכיל את הקבוצה הריקה, נקרא $\sg$־אלגברה (קרי: סיגמא אלגברה). }
	
	או במילים אחרות: אלגברה שעובדת גם לאיחודים בני מנייה (ולא רק סופיים). 
	
	יש סיגמה־אלגבראות מעוצמת הרצף, אך אין סיגמה־אלגברה (לא סופית) מעוצמה בת־מנייה. 
	
	לסיכום, הקבוצות המדידות ב־$\Og$ הוא $\sg$־אלגברה, והמידה החיצונית היא $\sg$־אדטיבית (כלומר היא אינווריאנטית לאיחוד קבוצות זרות). 
	
	\defi{עבור $A \subset \Og$ מדידה, מסמנים ב־$m(A) = m^*(A)$, וזו נקראת \textit{מידת לבג} של $A$. }
	
	זה כמובן מרחיב את ההגדרה הקודמת לקבוצות אלמנטריות. זה לא תלוי ב־$\Og$ (כי מידה חיצונית לא תלויה ב־$\Og$). 
	
	\textbf{דוגמאות: }
	\begin{enumerate}
		\item אם $A_1, A_2 \dots \subset \Og$ מדידות, אז גם $\bigcap_{n = 1}^{\infty}A_n$ מדיד (כי יש שימור של איחוד אינסופי ושל משלים, ומכללי דה מורגן, כלומר $\bigcap A_n = \cl{\bigcup A_n^{c}}^{c}$). 
		\item איחוד בן־מנייה של מלבנים אלמנטריים (או איחוד כל המלבנים הרציונליים שמוכלים בקבוצה מדידה). 
		\item כל קבוצה פתוחה $A \subset \Og$ מדידה. 
		\item כל קבוצה סגורה $A \subset \Og$ מדידה, זאת כי: 
		\[ \Og \setminus A = \underbrace{\Int \Og \setminus A}_{\ \!\text{פתוחה}} \cup \ B \]
		כאשר $B \subset \partial \Og$, ו־$\partial \Og$ זניחה וממונוטוניות וסגירות לאיחוד סיימנו. 
	\end{enumerate}
	
	\nthlec{2}{3}
	
	\subsubsection{הגדרת מידה מחוץ ל־$\Og$}
	כרגע, קובעים מלבן אלמנטרי $\Og \subset \R^{2}$ ועובדים בתוכו. נרצה להתרחב לכל $\R^{2}$. 
	
	מומלץ לקרוא על המאמר של Falconer מ־1986. 
	
	תופעה מעניינת: אין קבוצה מדידה שמכילה ''חצי מהממשיים בכל סקאלה``. הכוונה היא שלא קיימת $A \subset \R$ מדידה כך שלכל $I \subset \R$ מתקיים $m(A \cup I) = \frac{m(I)}{2}$. 
	
	כשנרצה להתרחב מחוץ ל־$\Og$, נרצה להגדיר גם מידה לקבוצות לא חסומות. 
	
	\textbf{הרעיון: }נרצף את $\R^{2}$ עם מלבנים אלמנטריים, לדוגמה באמצעות $[n, n +1) \times [m, m + 1)$ לכל $n, m \in \Z$. נסמן את הריבועים בריצוף ב־$\Og_1, \Og_2 \dots $ ומן הסתם מתקיים ש־$\R^{2} = \biguplus^{\infty}_{i = 1} \Og_i$. 
	
	\defi{\label{defi:mesuareR2}קבוצה $A \subset \R^{2}$ תקרא \textit{מדידה לבג} אם $\forall n \co A \cap \Og_n$ מדידה ב־$\Og_n$, ונגדיר את \textit{מידת לבג של $A$} להיות: 
	\[ m(A) \cod \bigcup_{n = 1}^{\infty}m(A \cap \Og_n) \]}
	
	\exe{נניח ש־$\Og_1, \Og_2$ מלבנים אלמנטריים, ונניח ש־$A \subset \Og_1 \cap \Og_2$. אז אם $A$ מדידה ביחס ל־$\Og_1$ היא גם מדידה ביחס ל־$\Og_2$. }
	\cola{הגדרת המדידות לא תלויה בבחירת הריצוף. }
	
	\clm{ההגדרה ב־\ref{defi:mesuareR2} מקיימת את התנאים הבאים: 
	\begin{enumerate}
		\item אוסף הקבוצות ב־$\R^{2}$ הוא $\sg$־אלגברה. 
		\item אם $A_1, A_2 \dots \subset \R^{2}$ מדידות וזרות, אז: 
		\[ m\cl{\bigcup_{n = 1}^{\infty}A_n\!} = \sum_{n = 1}^{\infty}m(A_n) \]
	\end{enumerate}}
	\begin{proof}
		\begin{enumerate}
			\item ברור ש־$\vn$ מדידה (כי $\vn \cap \Og_n = \vn$ מדידה ב־$\Og_n$). סגירות למשלים ברורה כי: 
			\[ \R^{2} \setminus A = \biguplus{\Og _n \setminus A} \]
			ו־$\Og_n \setminus A$ מדידה לכל $n$ (כי שווה ל־$\Og_n\setminus (\Og_n \cap A)$). ישנה סגירות לאיחוד בן־מנייה: 
			\[ \cl{\bigcup_{i = 1}^{\infty}A_i} \cap \Og_n = \bigcup_{i= 1}^{\infty}\underbrace{\cl{A_i \cap \Og _n}}_{\text{מדידה ב־$_n\Og$}} \]
			\item יהיו $A_1, A_2 \dots $ כמו קודם, ואז: 
			\[ \begin{WithArrows}
				&m\cl{\bigcup_{i =1}^{\infty}A_i} \Arrow[up]{מההגדרה}\\
				= &\sum_{n = 1}^{\infty}m\cl{\cl{\bigcup_{i =1}^{\infty}A_i} \cap \Og_n}  \\
				 = & \sum_{n = 1}^{\infty}m\cl{\bigcup_{i = 1}^{\infty}\cl{A_i \cap \Og_n}} \Arrow[]{$\sg$־אדטיביות ב־$_n\Og$} \\
				 = &\sum_{n = 1}^{\infty}\sum_{i= 1}^{\infty}m\cl{A_i \cap \Og_n} \Arrow[]{החלפת סדר סכימה} \\
				 = & \sum_{i= 1}^{\infty}\cl{\sum_{n = 1}^{\infty}m\cl{A_i \cap \Og_n}}
			\end{WithArrows} \]
		\end{enumerate}
	\end{proof}
	
	\clm{לכל $A \subset \R^{2}$ מדידה לבג, מתקיים $m(A) = m^*(A)$}
	\rmark{זכרו שמידה חיצונית מוגדרת לכל $\R^{2}$. ראה \ref{defi:outermeasure}. }
	\begin{proof}
		למה $m(A) \ge m^*(A)$? כי: 
		\[ m(A) = \sum_{n =1}^{\infty}m(A \cap \Og _n) = \sum_{n = 1}^{\infty}m^*(A \cap \Og_n) \ge m^*\cl{\bigcup_{n = 1}^{\infty}\cl{A \cap \Og _n}} = m^*(A) \]
		למה $m(A) \le m^*(A)$? מספיק להראות את השוויון $m^*(A) \ge \sum_{n = 1}^{N}m(A \cap \Og _n)$ לכל $N$, ואז לקחת גבול וסיימנו. נקבע מלבן אלמנטרי שמכיל את $\Og_1 \dots \Og_n$ (יש מספר סופי, ולכן יש מלבן אלמנטרי שמכיל את כולם). הקבוצה $A \cap \Og_i$ מדידה ב־$\Og$ לכל $i$ (תרגיל למעלה), לכן $A \cap \Og = \bigcup_{i = 1}^{N}\cl{A \cap \Og_i}$ מדידה ב־$\Og$, ו־: 
		\[ m^*(A) \ge m^*(A \cap \Og) m(A \cap \Og) \ge \sum_{i =1}^{N}m(A \cap \Og_i) \]
		כאשר הא''שים לעיל נכונים ממונוטוניות, מדידות ותת־אדטיביות. 
		והוכחנו את הנדרש. 
	\end{proof}
	
	בפרט, הגדרת מידת לבג לא תלויה בריצוף. 
	
	\textbf{שאלה: }נתונות $A, B \subset \R^{2}$ מדידות. האם $m(A \cup B) = m(A) + m(B)- m(A \cap B)$? 
	
	\textbf{תשובה: }זה לאו דווקא מוגדר היטב. ייתכן שהקבוצות אין־סופיות, ואי אפשר לחסר אינסוף מאינסוף. מה שאפשר להגיד הוא ש־$m(A \cup B) + m(A \cap B) = m(A) + m(B)$. 
	
	\textbf{שאלה: }איזה קבוצות מדידות יש ב־$\R^{2}$? הוכחנו שיש: 
	\begin{itemize}
		\item אלמנטריות
		\item פתוחות (כאיחוד בן מנייה של אלמנטריות) וסגורות
		\item קבוצות ממידה (חיצונית) אפס
		\item ועוד
	\end{itemize}
	
	שאלה מתבקשת היא כמה קבוצות מדידות יש ב־$\R^{d}$. 
	\clm{יש $2^{\con}$ קבוצות מדידות ב־$\R^{d}$. }\begin{proof}
		החסם העליון פשוט, כי $\sof{\ps(\R^{d})} = 2^{\con}$. נדרש למצוא רק חסם תחתון. נחפש קבוצה מעוצמת הרצף, שכל תת־הקבוצות שלה מדידות. מספיק למצוא קבוצה ממידה אפס, שעוצמתה רצף. ב־$d \ge 2$ ההוכחה פשוטה: ניקח פשוט קו ישר. ב־$d = 1$ נוכל להשתמש בקבוצת קנטור. 
	\end{proof}
	
	נתבונן בקבוצה החביבה הבאה: 
	\[ A_\eg = \bigcup_{n = 1}^{\infty}\cl{q_n - \frac{\eg}{2^{n}}, q + \frac{\eg}{2^{n}}} \]
	זו בבירור קבוצה פתוחה, $\Q \subset A_\eg$, ו־$A_\eg \subsetneq \R$ (כי מידתה לכל היותר $2\eg$ סופית מתת־סיגמה אדטיביות. ראינו כבר את זה בעבר). זו קבוצה פתוחה, צפופה ו''קטנה``. נוכל להקטין אותה אפילו עוד: 
	\[ A = \bigcap_{k = 1}^{\infty} A_{\frac{1}{k}} \]
	זו קבוצה אפילו יותר קטנה (כי $m(A) \le m(A_{\frac{1}{k}}) \le \frac{2}{k} \to 0$), צפופה (כי $\Q\subset A$), והיא גם מעוצמת הרצף (ממשפט הקטגוריה של בייר). הקבוצה $A$ היא קבוצת $G_\dg$. 
	\defi{תהא $A \subset X$ מרחב טופולוגי. $A$ תקרא $G_\dg$ אם קיימות $(O_i)_{i = 1}^{\infty}$ פתוחות כך ש־$A = \bigcap_{i = 1}^{\infty} O_i$, ו־$F_\sg$ אם קיימות $(F_i)_{i = 1}^{\infty}$ סגורות כך ש־$A = \bigcup_{i = 1}^{\infty}F_i$. }
	
	\exe{חיתוך של $\sg$־אלגבראות הוא $\sg$־אלגברה. }
	
	\defi{ה־$\sg$ של בורל, היא ה־$\sg$־אלגברה המינימלית שמכילה קבוצות פתוחות ב־$\R^{d}$. }
	\rmark{בלעז, Borel. }
	\rmark{ה־$\sg$־אלגברה של בורל מוגדרת ביחידות מכיוון שניתן לקחת את החיתוך של כל ה־$\sg$־אלגבראות ולקבל $\sg$־אלגברה מינימלית. }
	מסתבר שה־$\sg$־אלגברה של בורל קטנה יותר מה־$\sg$־אלגברה של לבג (קבוצת הקבוצות המדידות). 
	\clm{ה־$\sg$־אלגברה של בורל, היא הקבוצה המורכבת מהקבוצות הבאות: 
	\begin{itemize}
		\item קבוצות פתוחות
		\item קבוצות סגורות
		\item קבוצות $F_\sg$
		\item קבוצות $G_\dg$
	\end{itemize}}
%	באינדוקציה טרנספיניטית אפשר לייצר 

	נרצה לנסח קריטריון למדידות של קבוצות ב־$A \subset \R^{2}$ שלא מערב ריצוף ספציפי של $\R^{2}$. 
	
	\begin{Theorem}[תנאי המדידות של Lebesgue]
		תהא $A \subset \R^{2}$ קבוצה. אז $A$ מדידה אמ''מ לכל $\eg > 0$ קיימות $G$ פתוחה ו־$F$ סגורה כך ש־$F \subset A \subset G \land m(G \setminus F) < \eg$. 
	\end{Theorem}\begin{proof}
		\begin{itemize}
			\item[$\implies$] רוצים להוכיח ש־$A \cap \Og_n$ מדידה לכל $n$. נבחר $\Og \cod \Og_n$, ונסמן $A' = A \cap \Og$. רוצים להראות ש־$A'$ מדידה ב־$\Og$, כלומר שיש לה קירוב אלמנטרי טוב. יהי $\eg > 0$, ומתנאי המשפט נמצא $F \subset A \subset G$ מתאימות. נתבונן ב־$F' = F \cap |og, G' = G \cap \Og$. ברור ש־$F', G'$ מדידות ב־$\Og$ (כי $G' \setminus \ \partial \Og$ פתוחה ו־$F' \setminus \partial \Og $ סגורה, והשפה זניחה). 
			
			חיתוך של מדידות מדיד ב־$\R^{2}$ ולכן $G', F'$ מדידות ב־$\R^{2}$. לכן ל־$F', G'$ יש קירובים אלמנטריים $P_\eg, Q_\eg \subset \Og$ כך ש־$m^*(F' \trio P_\eg) ,m^*(G' \trio Q_\eg) < \eg$. נבחין ש־$P_\eg$ קירוב טוב: 
			\[ m^*(A' \trio P_\eg) \le m^*(\underbrace{A' \trio F'}_{A' \setminus F'}) + m^*(F' \trio P_\eg) \le m(A \setminus F) + \eg \le m(G \setminus F) + \eg \le 2\eg \]
			\item[$\impliedby$]נתון ש־$A$ מדידה, ולכל $\eg > 0$ נרצה למצוא $F, G$ המקיימים את התנאים. נגדיר (לצורך ההוכחה הזו) פנים של מלבן אלמנטרי־פתוח``. היה תרגיל לפיו מידה חיצונית אינווריאנטית ביחס ללקיחת פנים. 
			
			יהי $\eg > 0$. נבנה פתוחה $G$ כך ש־$m(G \setminus A) < \eg$. לכל $n \in \N$ ידוע ש־$m(A \cap \Og_n) = m^*(A \cap \Og_n)$ ולכן מהגדרת המידה החיצונית קיימים מלבנים אלמנטריים פתוחים $P_{n}^{(1)}, P_{n}^{(2)} \dots \subset \Og_n$ שקרובים לאינפימום כרצוננו ובפרט בקצב מעריכי (ישירות מהגדרת $\inf$ וממדידות), כלומר:
			\begin{align*}
				A \cap \Og = \bigcup_{i = 1}^{\infty}P_n^{(i)} && \sum_{i = 1}^{\infty}m(P_{n}^{(i)}) \le \frac{\eg}{2^{n}} + m(A \cap \Og_n)
			\end{align*}
			נשים לב ש־: 
			\[ \sum_{i = 1}^{\infty}m\cl{P_n^{(i)} \setminus A} \le \frac{\eg}{2^{n}} \]
			כי: 
			\[ m(P_{n}^{i} \setminus A) = m(P_n^{(i)}) - m(P_n^{(i)} \cap A) \]
			(הכל מתכנס בהחלט) דהיינו: 
			\[ \sum_{i = 1}^{\infty} \cl{m(P_n^{(i)}) - m(P_n^{(i)} \cap A)} = \frac{\eg}{2^{n}} + m(A \cap \Og_n) - \underbrace{\sum_{i = 1}^{\infty}m(P_n^{(i)} \cap A)}_{\mathclap{m\cl{\bigcup^{\infty}_{i = 1}(P_n^{(i)}\cap A)}}} = \frac{\eg}{2^{n}} + \cancel{m(A \cap \Og_n) - m(A \cap \Og_n)} \]
			(תנסו לצייר את זה, זה אמור להיות יחסית אינטואיטיבי). 
			עתה נגדיר: 
			\[ G \cod \bigcup_{n, i}P_n^{(i)} \]
			זו פתוחה (איחוד פתוחות), שמכילה את $A$ (כי לכל $n$ מתקיים $G \supset \bigcup_{i = 1} P_n^{(i)} \supset A \cap \Og_n$ ובאיחוד מוכלל על $n$ זה אכן מכיל את $A$). נחשב את $m(G \setminus A)$: 
			\[ m(G \setminus A) = m\!\cl{\bigcup_{n, i}P_n^{(i)}\setminus A} \le \sum_{n, i}m(P_n^{(i)}\setminus A) \le \sum_{n =1}^{\infty}\frac{\eg}{2^{n}} = \eg \]
			
			איך נמצא את $F$? נחזור על התהליך בעבור $\R^{2} \setminus A$, ונקב ל$G'$ פתוחה כך ש־$\R^{2} \setminus \subset G'$ וכן $m(G' \setminus (\R^{2} \setminus A))) < \eg$ ומכאן קל לראות ש־$F = \R^{2} \setminus G'$ מקיימת את הנדרש. 
			
		\end{itemize}
	\end{proof}
	
%	\exe{התנאים הבאים שקולים: 
%	\begin{itemize}
%		\item לכל $\eg$ קיימת $F_\eg$ סגורה שמקרבת טוב
%		\item קיימת $G_\eg$ פתוחה שמקרבת טוב
%		\item $A$ מדידה
%		\end{itemize}}

%TODO check above
	
	\cola{תהא $A \subset \R^{2}$ מדידה. אז: 
	\begin{align*}
		m(A) &=\inf\ccb{m(G) \vert A \subset G \land G \ \text{\en{open}}} \\
		&= \inf\ccb{m(F) \vert F \subset A \land F \ \text{\en{closed}}}
		\end{align*}}\begin{proof}
		ברור ש־$m(A) \le \inf \ccb{\cdots}$. מהצד השני, לכל $\eg > 0$, $G$ פתוחה ו־$F$ סגורה, נקבל ש־: 
		\[ m(G \setminus A) \le m(G \setminus P) < \eg \]
		ולכן $m(G) = m(A) + m(G \setminus A) \le m(A) + \eg$ וסיימנו. הסופרמום באופן דומה. 
	\end{proof}
	
	
	\exe{לכל $A \subset \R^{d}$ מדידה, מתקיים ש־: 
	\[ m(A) = \sup\ccb{m(K) \mid K \subset A \land K \ \text{\en{compact}}} \]}
	
	\cola{לכל מדידה $A \subset \R^{2}$ קיימת $S \subset \R^{2}$ שמכילה את $A$ והיא $G_\dg$, כך ש־$m(S \setminus A)= 0$. }\begin{proof}
		לכל $n$ יש $G_n$ פתוחה שמכילה את $A$ ו־$m(G_n \setminus A) \le \frac{1}{n}$, ונגדיר: 
		\[ S = \bigcap_{i = 1}^{\infty}G_n \]
		מהגדרה היא $G_n$ ו־: 
		\[ m(S \setminus A) \le m(G_n \setminus A) \le \frac{1}{n} \toinf 0 \]
	\end{proof}
	
	בכך סיימנו את בניית מידת לבג ב־$\R^{2}$. 
	
	נבחין במספר תכונות של מידת לבג שמצאנו: 
	\clm{נניח ש־$A \subset \R^{d}$ מדידה. אז: 
	\begin{itemize}
		\item אם $x \in \R^{d}$, אז $A + x$ מדידה. 
		\item אם $\lg \in \R$, אז $m(\lg A) =\sof{\lg}^{d} m(A)$ מדידה. 
		\item אם $T \co \R^{d} \to \R^{d}$ לינארית, אז: 
		\[ m(T(A)) = \sof{\det T} m(A) \]
		(הוכחה של זה בחדו''א 3)q
		\end{itemize}}
	
	
	\section[מידת לבג למרחב מידה]{\en{Lebegue Measure for a Mesuare Space}}
	
	
	
	
	
	
	
	
	
	
	
	\ndoc
\end{document}
